%-------------------------
% Academic CV - Xin Sun
% Based on Jake's Resume template (MIT License)
% Compiled with XeLaTeX for CJK support
%-------------------------

\documentclass[letterpaper,11pt]{article}

\usepackage[UTF8]{ctex}
\usepackage{latexsym}
\usepackage[empty]{fullpage}
\usepackage{titlesec}
\usepackage{marvosym}
\usepackage[usenames,dvipsnames]{color}
\usepackage{verbatim}
\usepackage{enumitem}
\usepackage[hidelinks]{hyperref}
\usepackage{fancyhdr}
\usepackage{tabularx}
\usepackage{fontawesome5}
\usepackage{xcolor}

\pagestyle{fancy}
\fancyhf{}
\fancyfoot{}
\renewcommand{\headrulewidth}{0pt}
\renewcommand{\footrulewidth}{0pt}

% Adjust margins
\addtolength{\oddsidemargin}{-0.6in}
\addtolength{\evensidemargin}{-0.6in}
\addtolength{\textwidth}{1.2in}
\addtolength{\topmargin}{-0.7in}
\addtolength{\textheight}{1.4in}

\urlstyle{same}

\raggedbottom
\raggedright
\setlength{\tabcolsep}{0in}

% Define colors
\definecolor{primary}{RGB}{0, 79, 144}
\definecolor{secondary}{RGB}{100, 100, 100}

% Section formatting
\titleformat{\section}{
  \vspace{-8pt}\scshape\raggedright\large\bfseries\color{primary}
}{}{0em}{}[\color{primary}\titlerule \vspace{-6pt}]

% Custom commands
\newcommand{\resumeItem}[1]{
  \item\small{#1 \vspace{-2pt}}
}

\newcommand{\resumeSubheading}[4]{
  \vspace{-2pt}\item
    \begin{tabular*}{0.97\textwidth}[t]{l@{\extracolsep{\fill}}r}
      \textbf{#1} & #2 \\
      \textit{\small#3} & \textit{\small #4} \\
    \end{tabular*}\vspace{-7pt}
}

\newcommand{\resumeSubHeadingListStart}{\begin{itemize}[leftmargin=0.15in, label={}]}
\newcommand{\resumeSubHeadingListEnd}{\end{itemize}}
\newcommand{\resumeItemListStart}{\begin{itemize}[leftmargin=0.2in]}
\newcommand{\resumeItemListEnd}{\end{itemize}\vspace{-5pt}}

% Publication entry
\newcommand{\pubentry}[4]{
  \item \textbf{#1} \\
  {\small #2} \\
  {\small \textit{#3}\hfill #4}
  \vspace{0pt}
}

%-------------------------------------------
\begin{document}

%----------HEADING----------
\noindent
\begin{tabular*}{\textwidth}{l@{\extracolsep{\fill}}r}
  {\Huge \textbf{孙鑫}} \; {\large Xin Sun} &
  \small \faPhone\ (+86) 193-8968-7089 \;\; \faEnvelope\ \href{mailto:xin.sun@cripac.ia.ac.cn}{xin.sun@cripac.ia.ac.cn} \\[4pt]
  {\small 中国科学技术大学 \; / \; 中国科学院自动化研究所} &
  \small \faGraduationCap\ \href{https://scholar.google.com/citations?user=x2aQKE4AAAAJ}{Google Scholar} \;\; \faGithub\ \href{https://github.com/sunxin000}{GitHub}
\end{tabular*}
\vspace{-6pt}
{\color{primary}\rule{\textwidth}{1.5pt}}

%----------RESEARCH INTERESTS----------
\section{研究方向}
我的研究聚焦于\textbf{检索增强生成 (RAG)} 与\textbf{大模型智能体训练}，涵盖两个核心方向：
\textbf{(1) RAG \& Deep Search}——构建可靠、自适应的检索增强系统，使模型知道何时、检索什么、如何检索，迈向自主深度研究智能体；
\textbf{(2) LLM Agent Training}——通过强化学习与后训练赋予大模型智能体多轮探索、工具调用、结构化推理及环境反馈学习等能力。

%----------EDUCATION----------
\section{教育背景}
  \resumeSubHeadingListStart
    \resumeSubheading
      {中国科学技术大学}{合肥, 中国}
      {控制科学与工程 \quad 在读博士（中国科学院自动化研究所联合培养）}{2022 -- 2027（预计）}
      \resumeItemListStart
        \resumeItem{中国科学院自动化研究所模式识别国家重点实验室 (NLPR)}
        \resumeItem{导师：王亮 教授（IEEE Fellow）、刘强 副研究员}
        \resumeItem{荣誉：研究生国家奖学金（2024）、中国科学技术大学一等学业奖学金（2022--2025, 连续三年）}
      \resumeItemListEnd

    \resumeSubheading
      {上海交通大学}{上海, 中国}
      {计算机科学与技术 \quad 工学学士}{2018 -- 2022}
      \resumeItemListStart
        \resumeItem{荣誉：国家励志奖学金（2019--2022）、上海交通大学学业奖学金}
      \resumeItemListEnd
  \resumeSubHeadingListEnd

%----------PUBLICATIONS----------
\section{论文发表}

\textbf{\color{primary} $\blacktriangleright$ \quad RAG, Deep Search \& Agentic RL（核心研究方向）}
\begin{enumerate}[leftmargin=0.3in, label={[\arabic*]}]

  \pubentry{KBQA-R1: Reinforcing Large Language Models for Knowledge Base Question Answering}
  {\textbf{\underline{Xin Sun}}, Zhongqi Chen, Xing Zheng, Qiang Liu, Shu Wu, Bowen Song, Zilei Wang, Weiqiang Wang, Liang Wang}
  {Submitted to ICML 2026 \textnormal{(CCF-A, Scores: 4/4/4/5, 全正分)}}{}

  \pubentry{Divide-Then-Align: Honest Alignment based on the Knowledge Boundary of RAG}
  {\textbf{\underline{Xin Sun}}, Jianan Xie, Zhongqi Chen, Qiang Liu, Shu Wu, Yuehe Chen, Bowen Song, Zilei Wang, Weiqiang Wang, Liang Wang}
  {ACL 2025, Main Conference (CCF-A)}{}

  \pubentry{Predict the Retrieval! Test Time Adaptation for Retrieval Augmented Generation}
  {\textbf{\underline{Xin Sun}}, Zhongqi Chen, Qiang Liu, Shu Wu, Bowen Song, Weiqiang Wang, Zilei Wang, Liang Wang}
  {ICASSP 2026 (CCF-B)}{}

  \pubentry{ToolWeaver: Weaving Collaborative Semantics for Scalable Tool Use in Large Language Models}
  {Bowen Fang, Qiang Liu, Yesheng Liu, Jiabing Yang, \textbf{\underline{Xin Sun}}, Shu Wu, Baole Wei, Liang Wang}
  {ICLR 2026 (CAAI-A)}{}

  \pubentry{Multimodal Adaptive Retrieval Augmented Generation through Internal Representation Learning}
  {Ruoshuang Du, \textbf{\underline{Xin Sun}}, Qiang Liu, Bowen Song, Zhongqi Chen, Weiqiang Wang, Liang Wang}
  {IJCNN 2026 (CCF-C)}{}

\end{enumerate}

\vspace{2pt}
\textbf{\color{primary} $\blacktriangleright$ \quad 其他研究工作：图学习 \& 可信机器学习}
\begin{enumerate}[leftmargin=0.3in, label={[\arabic*]}, resume]

  \pubentry{DIVE: Subgraph Disagreement for Graph Out-of-Distribution Generalization}
  {\textbf{\underline{Xin Sun}}, Liang Wang, Qiang Liu, Shu Wu, Zilei Wang, Liang Wang}
  {KDD 2024 (CCF-A)}{}

  \pubentry{Noise-Robust Semi-Supervised Learning for Distant Supervision Relation Extraction}
  {\textbf{\underline{Xin Sun}}, Qiang Liu, Shu Wu, Zilei Wang, Liang Wang}
  {EMNLP 2023, Findings (CCF-B)}{}

  \pubentry{Uncertainty Calibration for Counterfactual Propensity Estimation in Recommendation}
  {Wenbo Hu*, \textbf{\underline{Xin Sun}}*, Qiang Liu, Liang Wang}
  {IEEE TKDE (CCF-A, SCI Q1)}{}

  \pubentry{Pin-Tuning: Parameter-Efficient In-Context Tuning for Few-shot Molecular Property Prediction}
  {Liang Wang, Qiang Liu, Shaozhen Liu, \textbf{\underline{Xin Sun}}, Shu Wu, Liang Wang}
  {NeurIPS 2024 (CCF-A)}{}

  \pubentry{GSLB: The Graph Structure Learning Benchmark}
  {Zhixun Li, Liang Wang, \textbf{\underline{Xin Sun}}, Yifan Luo, Yanqiao Zhu, Dingshuo Chen, Yingtao Luo, Xiangxin Zhou, Qiang Liu, Shu Wu, Liang Wang, Jeffrey Xu Yu}
  {NeurIPS 2023 (CCF-A)}{}

\end{enumerate}

%----------PROJECTS----------
\section{开源项目}
  \resumeSubHeadingListStart
    \item
    \begin{tabular*}{0.97\textwidth}[t]{l@{\extracolsep{\fill}}r}
      \textbf{GSLB: 图结构学习评测基准} & \href{https://github.com/GSL-Benchmark/GSLB}{\faGithub\ GitHub \quad \faStar\ 103} \\
    \end{tabular*}\vspace{-4pt}
    \resumeItemListStart
      \resumeItem{面向图结构学习的开源评测工具库，集成多种结构学习算法，收录大量经典模型复现代码}
    \resumeItemListEnd
  \resumeSubHeadingListEnd

%----------EXPERIENCE----------
\section{实习经历}
  \resumeSubHeadingListStart
    \resumeSubheading
      {蚂蚁集团}{杭州, 中国}
      {机器智能部门 -- 算法研究实习生}{2024.11 -- 2026.04}
      \resumeItemListStart
        \resumeItem{RAG 核心方向研究：主导完成多项检索增强生成相关工作，包括 RAG 可信对齐 (ACL 2025)、多模态自适应 RAG (IJCNN 2026)、RAG 测试时自适应 (ICASSP 2026)、知识库问答强化学习 KBQA-R1 (投稿 ICML 2026) 等}
      \resumeItemListEnd

    \resumeSubheading
      {字节跳动}{北京, 中国}
      {AML -- 机器学习工程师实习生}{2022.02 -- 2022.08}
      \resumeItemListStart
        \resumeItem{参与多变量时间序列自动机器学习平台的开发}
      \resumeItemListEnd
  \resumeSubHeadingListEnd

%----------SERVICE----------
\section{学术服务}
\begin{itemize}[leftmargin=0.15in, label={}, nosep]
  \item \textbf{会议审稿人：} ICML 2026, NeurIPS 2025, ACL 2025, EMNLP 2025, KDD 2024
\end{itemize}

%-------------------------------------------
\end{document}
